% !TeX root = ../beamer.tex

\section{Data and citations}

\subsection{Table}
\begin{frame}{Table}
    \begin{table}
        \begin{tabular}{l r r}
            \textbf{Group} & \textbf{Mean} & \textbf{Deviation} \\
            \headrule
            Control             & 0.562               & 0.0123              \\ \hline
            Low dose            & 0.910               & 0.0456              \\ \hline
            High dose           & 0.296               & 0.0089              \\
        \end{tabular}
        \caption{Three groups of the experiment}
    \end{table}
    \source{Verdi (2021), Table 2}
\end{frame}

\subsection{Table: code and result}
\begin{frame}{Table: code and result}
\begin{columns}[T,onlytextwidth]
\begin{column}{0.5\textwidth}
\step{1. Header and rule}
\begin{codeonly}
\textbf{Group} & \textbf{Mean} \\
\headrule
\end{codeonly}
\smallskip
\step{2. Rows}
\begin{codeonly}
Control & 0.562 \\ \hline
Treated & 0.910 \\
\end{codeonly}
\smallskip
\step{3. Source}
\begin{codeonly}
\source{Verdi (2021)}
\end{codeonly}
\end{column}
\begin{column}{0.45\textwidth}
\stepoffset
\begin{table}
  \begin{tabular}{l r}
    \textbf{Group} & \textbf{Mean} \\
    \headrule
    Control & 0.562 \\ \hline
    Treated & 0.910 \\
  \end{tabular}
  \caption{Group means}
\end{table}
\source{Verdi (2021)}
\end{column}
\end{columns}
\end{frame}

\subsection{Figure: transparent}
\begin{frame}{Figure: transparent}
\begin{columns}[T,onlytextwidth]
\begin{column}{0.5\textwidth}
\step{1. A figure with alpha}
\begin{codeonly}
\begin{figure}
  \includegraphics[height=4cm]
    {graph.png}
  \caption{Caption}
\end{figure}
\end{codeonly}
\smallskip
\step{2. Name the source}
\begin{codeonly}
\source{AVDLCZ, CC0}
\end{codeonly}
\end{column}
\begin{column}{0.45\textwidth}
\stepoffset
\begin{figure}
    \includegraphics[height=0.5\textheight]{graph.png}
    \caption{A PNG with a transparent background sits on the slide as it is}
\end{figure}
\source{AVDLCZ, CC0}
\end{column}
\end{columns}
\end{frame}

\subsection{Figure: opaque}
\begin{frame}{Figure: opaque}
\begin{columns}[T,onlytextwidth]
\begin{column}{0.5\textwidth}
\step{1. A photo, corners rounded}
\begin{codeonly}
\begin{figure}
  \roundedgraphics
    [height=4cm]{photo.jpg}
  \caption{Caption}
\end{figure}
\end{codeonly}
\smallskip
\step{2. Name the source}
\begin{codeonly}
\source{NASA / B. Anders}
\end{codeonly}
\end{column}
\begin{column}{0.45\textwidth}
\stepoffset
\begin{figure}
    \roundedgraphics[height=0.5\textheight]{earthrise.jpg}
    \caption{A photo keeps its own background, with the corners slightly rounded}
\end{figure}
\source{NASA / B. Anders}
\end{column}
\end{columns}
\end{frame}

\subsection{Citation}
\begin{frame}{Citation}
    \verb|\cite| gives the number of the entry in the list of references.

    One work is cited like this \cite{p1}, two at once like this \cite{p1,p2}.
\end{frame}

\subsection{Citation: code and result}
\begin{frame}{Citation: code and result}
\begin{columns}[T,onlytextwidth]
\begin{column}{0.5\textwidth}
\step{1. Cite}
\begin{codeonly}
\cite{p1}
\cite{p1,p2}
\end{codeonly}
\smallskip
\step{2. Name the sources}
\begin{codeonly}
\source{Verdi (2021)}
\source{Rossi (2019)}
\end{codeonly}
\smallskip
\step{3. Lighten a remark}
\begin{codeonly}
\fade{remark}
\end{codeonly}
\end{column}
\begin{column}{0.45\textwidth}
\stepoffset
One source \cite{p1}, two at once \cite{p1,p2}, and a faded \fade{remark}.
\source{Verdi (2021)}\source{Rossi (2019)}
\end{column}
\end{columns}
\end{frame}
